\section{Limitations of Back-Side Power Delivery Networks}

The improvements in power integrity and routing efficiency offered by BSPDNs come at a cost. Moving the power grid to the backside of the wafer introduces thermal, self-heating, and AC impedance challenges that are not present in conventional front-side architectures.

The most pressing issue is thermal performance. In FSPDNs, the full-thickness silicon substrate spreads heat laterally from the active devices toward the topside heat sink. BSPDNs require thinning the wafer to form nano through-silicon via (nTSV) connections, which removes much of this lateral heat-spreading path. The bonding interface and back-end-of-line (BEOL) layers between the devices and the heat sink further increase thermal resistance. Modelling of CPU hotspot regions has shown that BSPDN configurations reach temperatures approximately 45\% higher than equivalent FSPDN designs\cite{10925422}. Several mitigation strategies have been investigated, including high thermal conductivity bonding materials, stacked BEOL via configurations that create direct vertical thermal paths, and embedded microchannel cooling etched into the backside metal layers\cite{10925422}.

At the device level, the thinned substrate also worsens self-heating effects (SHE). With less silicon to absorb and spread the heat generated by switching transistors, localised temperature rises become more severe. An evaluation of SHE in SRAM macros and circuits at the 3\,nm node found that BSPDN designs improve power, performance, and area, but the resulting thermal behaviour requires power-thermal co-design to avoid reliability degradation\cite{10529350}. The study further showed that the choice of buried power rail material and the number of backside metal layers both have a measurable impact on device temperature, indicating that thermal optimisation must be considered alongside electrical design choices from the earliest stages of the design flow\cite{10529350}. Electrical gains therefore cannot be evaluated in isolation from their thermal consequences.

A separate concern relates to the AC impedance profile of the backside power grid. Because the DC resistance of BSPDNs is much lower than that of FSPDNs, the damping in the power delivery network is reduced. As a result, impedance resonance peaks can shift into frequency ranges that coincide with circuit switching activity, which risks supply voltage fluctuations during transient load changes\cite{8993617}. Studies on the Arm Cortex-A53 benchmark showed that this resonance effect is particularly relevant in multi-core designs, where on-chip decoupling capacitors are needed to suppress ringing\cite{8993617}. Careful impedance profiling and decoupling capacitor placement during the design phase are therefore essential when adopting a backside power architecture.

